% ============================================================================
% DÉFINITION DU COURS 13 : MACHINES ASYNCHRONES ET SYNCHRONES
% ============================================================================

\titlehead{Thomas Lusseau - SII en ATS}
\title[MAS et MS]{Machines asynchrones et synchrones}
\author[TL]{Thomas Lusseau}
\date{\today}

\TLCourseCoverSetup
  {\repStyle/FondPoly.png}
  {\repStyle/FondPoly.png}
  {Cycle 7 : Analyser, modéliser, expérimenter et résoudre la distribution\\et la conversion d'énergie en alternatif}
  {Machines asynchrones\\et synchrones}
  {Lycée Robert Doisneau - ATS}
  {Thomas Lusseau}

\TLCourseCoverContents{%
  \TLCourseContentsLine{13.1}{Distribution triphasée et complexes}{sec:masms-1}
  \TLCourseContentsLine{13.2}{Onduleur triphasé}{sec:masms-2}
  \TLCourseContentsLine{13.3}{Généralités sur les machines alternatives}{sec:masms-3}
  \TLCourseContentsLine{13.4}{Machine asynchrone}{sec:masms-4}
  \TLCourseContentsLine{13.5}{Machine synchrone}{sec:masms-5}
  \TLCourseContentsLine{13.6}{Synthèse et démarche d'étude}{sec:masms-6}
  \TLCourseContentsLine{13.7}{Sources}{sec:masms-7}
  \TLCourseContentsLine{13.8}{Exercices du chapitre}{sec:masms-8}
}

\TLCourseFooter
\TLCourseCover
\markboth{Machines asynchrones et synchrones}{Machines asynchrones et synchrones}

\livrettrue
\collefalse

\TLStartCourseChapter{13}{Machines asynchrones et synchrones}{chap:mas-ms}

% ============================================================================
% COURS 13 — MACHINES ASYNCHRONES ET SYNCHRONES
% Partie A : triphasé, onduleur et généralités
% Source : 17-MAS_MS.docx
% ============================================================================

\newcommand{\TLMACheck}{\ensuremath{\square}}
\newcommand{\TLMAImage}[2][.78\linewidth]{%
  \IfFileExists{pandoc/media/#2}{%
    \begin{center}\includegraphics[width=#1]{pandoc/media/#2}\end{center}%
  }{}%
}
\newcommand{\TLMAInlineImage}[2][.30\linewidth]{%
  \IfFileExists{pandoc/media/#2}{\includegraphics[width=#1]{pandoc/media/#2}}{}%
}

\begin{TLMachineMethod}{Objectifs du cours}
Donner les démarches et les outils nécessaires à l'étude des machines alternatives : réseau triphasé, onduleur, machine asynchrone, machine synchrone et commande à vitesse variable.
\end{TLMachineMethod}

\begin{center}
\small
\begin{tabularx}{\linewidth}{>{\bfseries}p{2.0cm}X c}
\toprule
Je connais & Compétences et connaissances attendues & \\
\midrule
& La relation entre tension simple et tension composée et la définition d'un système triphasé équilibré. & \TLMACheck\\
& Les représentations vectorielle et complexe d'une grandeur sinusoïdale ainsi que les impédances de base. & \TLMACheck\\
& Les puissances active, réactive et apparente et le théorème de Boucherot. & \TLMACheck\\
& Les couplages étoile et triangle d'une machine triphasée. & \TLMACheck\\
& La vitesse de synchronisme, le glissement, le bilan de puissance et le couple d'une MAS. & \TLMACheck\\
& Le modèle et le diagramme de Behn--Eschenburg d'une machine synchrone. & \TLMACheck\\
& Le principe d'une commande $V/f$ constante et de l'autopilotage d'une MS. & \TLMACheck\\
\midrule
Je sais & Déterminer les grandeurs de ligne, les puissances et le facteur de puissance d'un réseau triphasé. & \TLMACheck\\
& Choisir le couplage d'une machine à partir de sa plaque signalétique et du réseau. & \TLMACheck\\
& Déterminer vitesse de synchronisme, glissement, rendement et couple électromagnétique. & \TLMACheck\\
& Tracer et exploiter un diagramme vectoriel et une caractéristique couple--vitesse. & \TLMACheck\\
& Vérifier une contrainte thermique et sélectionner une loi de variation de vitesse. & \TLMACheck\\
\bottomrule
\end{tabularx}
\end{center}

\section{Distribution en triphasé et utilisation des complexes}
\label{sec:masms-1}

\subsection{Tensions simples et composées}

Une installation triphasée comporte trois conducteurs de ligne identiques, appelés \textbf{phases}, et éventuellement un quatrième conducteur appelé \textbf{neutre}. Un système triphasé équilibré est formé de trois grandeurs sinusoïdales de même valeur efficace $V$, de même fréquence $f$ et déphasées de $120^\circ$.

Pour le système direct :
\[
\begin{aligned}
v_1(t)&=V\sqrt2\sin(\omega t),\\
v_2(t)&=V\sqrt2\sin\left(\omega t-\frac{2\pi}{3}\right),\\
v_3(t)&=V\sqrt2\sin\left(\omega t-\frac{4\pi}{3}\right),
\end{aligned}
\qquad \omega=2\pi f.
\]

Le système est équilibré lorsque :
\[
\boxed{v_1(t)+v_2(t)+v_3(t)=0.}
\]

La \textbf{tension simple} $v_i$ est mesurée entre la phase $i$ et le neutre. Sa valeur efficace est notée $V$. La \textbf{tension composée} $u_{ij}=v_i-v_j$ est mesurée entre deux phases et sa valeur efficace est notée $U$.

\begin{center}\TLTriphasageSimpleCompose\end{center}

Pour un réseau équilibré :
\[
\boxed{U=\sqrt3\,V.}
\]
Ainsi, dans un réseau noté $230/400\ \mathrm V$, la tension simple vaut $230\ \mathrm V$ et la tension composée $400\ \mathrm V$.

\TLMAImage[.56\linewidth]{image5.png}
\TLMAImage[.50\linewidth]{image7.png}

\begin{TLMachineExercise}{Tensions d'un réseau triphasé}
Un réseau est annoncé $230/400\ \mathrm V$. Déterminer les valeurs efficaces des tensions phase--neutre et phase--phase. Représenter les trois tensions simples puis la tension composée $\underline U_{12}$.
\end{TLMachineExercise}

\subsection{Représentation vectorielle et complexe}

Une grandeur sinusoïdale
\[
s(t)=S\sqrt2\sin(\omega t+\varphi)
\]
est représentée par le complexe :
\[
\boxed{\underline S=S\,e^{j\varphi}.}
\]
La norme est la valeur efficace, directement mesurable, et l'argument est la phase par rapport à une grandeur choisie comme référence.

\begin{center}\TLTriphasageVectoriel\end{center}

La représentation de Fresnel transforme les équations différentielles en équations algébriques. Pour les dipôles de base :
\[
\underline Z_R=R,
\qquad
\underline Z_L=jL\omega,
\qquad
\underline Z_C=\frac{1}{jC\omega}.
\]
Les associations série et parallèle s'étudient respectivement par addition des impédances et des admittances.

\TLMAImage[.48\linewidth]{image11.png}
\TLMAImage[.38\linewidth]{image14.png}

\subsection{Puissances active, réactive et apparente}

Pour un récepteur monophasé :
\[
P=UI\cos\varphi,
\qquad
Q=UI\sin\varphi,
\qquad
S=UI,
\qquad
S^2=P^2+Q^2.
\]

Pour un récepteur triphasé équilibré, avec $U$ tension composée et $I$ courant de ligne :
\[
\boxed{P=\sqrt3\,UI\cos\varphi,}
\qquad
\boxed{Q=\sqrt3\,UI\sin\varphi,}
\qquad
\boxed{S=\sqrt3\,UI.}
\]

Le facteur de puissance vaut :
\[
f_p=\frac{P}{S}.
\]
Lorsque les tensions et courants sont sinusoïdaux, $f_p=\cos\varphi$.

\begin{center}
\small
\begin{tabularx}{\linewidth}{>{\bfseries}l c c X}
\toprule
Comportement & Déphasage & Puissance réactive & Interprétation \\
\midrule
Résistif & $\varphi=0$ & $Q=0$ & courant en phase avec la tension.\\
Inductif & $0<\varphi\leq\pi/2$ & $Q>0$ & courant en retard.\\
Capacitif & $-\pi/2\leq\varphi<0$ & $Q<0$ & courant en avance.\\
\bottomrule
\end{tabularx}
\end{center}

Pour les dipôles idéaux :
\[
\begin{array}{c|c|c}
\text{résistance} & \text{inductance} & \text{condensateur}\\ \hline
P=RI^2=U^2/R & P=0 & P=0\\
Q=0 & Q=L\omega I^2=U^2/(L\omega) & Q=-C\omega U^2=-I^2/(C\omega)
\end{array}
\]

Le \textbf{théorème de Boucherot} indique que, pour des grandeurs sinusoïdales de même fréquence, les puissances actives et réactives totales sont les sommes algébriques des puissances de chaque récepteur.

\begin{TLMachineWarning}{Machine synchrone et puissance réactive}
Selon son excitation, une machine synchrone peut présenter un comportement inductif ou capacitif. Surexcitée, elle peut fournir de la puissance réactive et être utilisée comme compensateur synchrone.
\end{TLMachineWarning}

\begin{TLMachineExercise}{Puissances en triphasé}
Un réseau $230/400\ \mathrm V$, $50\ \mathrm{Hz}$ alimente un récepteur triphasé constitué par phase d'une résistance $R=20\ \Omega$ et d'une inductance $L=1\ \mathrm H$. Déterminer $\underline Z$, le courant de ligne, $P$, $Q$, $S$ et le facteur de puissance.
\end{TLMachineExercise}

\section{Onduleur triphasé}
\label{sec:masms-2}

Pour modifier la vitesse d'une MAS ou d'une MS, il faut agir sur la fréquence d'alimentation. L'\textbf{onduleur triphasé} recrée un système triphasé à partir d'une tension continue, appelée \textbf{bus continu}, fournie par une batterie ou par un redresseur.

\begin{center}\TLOnduleurTriphas\end{center}

Il comporte trois bras. Dans chaque bras, les deux interrupteurs sont commandés de manière complémentaire afin d'éviter le court-circuit du bus continu. La commande fixe la fréquence du fondamental et sa valeur efficace.

\TLMAImage[.70\linewidth]{image17.png}
\TLMAImage[.72\linewidth]{image18.png}

Une commande en créneaux produit des harmoniques de basse fréquence. Une modulation de largeur d'impulsions, ou \textbf{MLI}, repousse les premiers harmoniques autour de la fréquence de porteuse. Les inductances de la machine jouent alors le rôle d'un filtre passe-bas et le courant reste presque sinusoïdal.

\begin{center}
\begin{tabular}{cc}
\TLMAInlineImage[.42\linewidth]{image21.png} & \TLMAInlineImage[.42\linewidth]{image22.png}
\end{tabular}
\end{center}
\TLMAImage[.72\linewidth]{image23.png}

La puissance active moyenne est transportée par le fondamental de la tension et du courant. Pour une phase :
\[
P_\varphi=V_1 I_1\cos\varphi_1,
\]
et pour les trois phases équilibrées :
\[
P=3V_1I_1\cos\varphi_1.
\]

\begin{TLMachineExercise}{Étude d'un onduleur}
À partir des formes d'onde de la tension $v_{AN}$ et du courant $i_a$ :
\begin{enumerate}
  \item identifier un bras d'onduleur et justifier la complémentarité de ses commandes ;
  \item déterminer la fréquence du fondamental et la valeur efficace de sa première harmonique ;
  \item expliquer pourquoi le courant est presque sinusoïdal ;
  \item déterminer $P$ et $Q$ à partir du fondamental de la tension et du courant.
\end{enumerate}
\end{TLMachineExercise}

\section{Généralités sur les machines alternatives}
\label{sec:masms-3}

\subsection{Constitution d'une machine triphasée}

Une machine alternative triphasée, asynchrone ou synchrone, comporte :
\begin{itemize}
  \item un \textbf{stator}, partie fixe, portant trois enroulements décalés de $120^\circ$ électriques ;
  \item un \textbf{rotor}, partie tournante : cage d'écureuil ou bobinage pour une MAS, aimants permanents ou enroulement d'excitation pour une MS ;
  \item une \textbf{plaque à bornes} à six bornes permettant le couplage des enroulements ;
  \item une \textbf{plaque signalétique} donnant les conditions nominales d'utilisation.
\end{itemize}

\TLMAImage[.42\linewidth]{image28.jpeg}

\subsection{Création du champ tournant}

Les trois enroulements statoriques, décalés dans l'espace de $120^\circ$ et alimentés par un système triphasé équilibré, créent un champ magnétique tournant de norme pratiquement constante.

\begin{center}\TLChampTournant\end{center}

\paragraph{Théorème de Leblanc.}
Un courant sinusoïdal parcourant une bobine fixe peut être remplacé par deux champs tournants de même amplitude, tournant en sens opposés à la vitesse $\Omega=\omega/p$.

\paragraph{Théorème de Ferraris.}
Trois courants triphasés équilibrés parcourant trois bobines décalées de $120^\circ$ créent un champ tournant unique, de vitesse constante :
\[
\boxed{\Omega_s=\frac{\omega}{p}.}
\]

Pour une MS, le rotor s'accroche au champ et tourne exactement à sa vitesse. Pour une MAS, les courants rotoriques induits produisent un couple de poursuite : le rotor tourne à une vitesse différente de celle du champ.

\TLMAImage[.65\linewidth]{image30.png}

\subsection{Vitesse de synchronisme}

Avec $p$ paires de pôles et une alimentation de fréquence $f$ :
\[
\boxed{\Omega_s=\frac{2\pi f}{p}\ \mathrm{rad\,s^{-1}},}
\qquad
\boxed{n_s=\frac{f}{p}\ \mathrm{tr\,s^{-1}},}
\qquad
\boxed{N_s=\frac{60f}{p}\ \mathrm{tr\,min^{-1}}.}
\]

\begin{TLMachineExercise}{Pôles et vitesse de synchronisme}
Déterminer :
\begin{enumerate}
  \item la vitesse de synchronisme d'une machine tétrapolaire alimentée à $50\ \mathrm{Hz}$ ;
  \item celle d'une machine bipolaire alimentée à $100\ \mathrm{Hz}$ ;
  \item le nombre de paires de pôles d'une MS tournant à $2400\ \mathrm{tr\,min^{-1}}$ sous $80\ \mathrm{Hz}$ ;
  \item le nombre de paires de pôles d'une MAS tournant à $950\ \mathrm{tr\,min^{-1}}$ sous $50\ \mathrm{Hz}$.
\end{enumerate}
\end{TLMachineExercise}

\subsection{Plaque signalétique et couplages}

La plaque signalétique précise notamment l'indice de protection, les tensions et courants nominaux, la fréquence, la vitesse nominale, le facteur de puissance et la puissance utile mécanique.

La plus petite tension portée sur la plaque correspond à la tension nominale d'un enroulement. Le couplage doit donc être choisi pour appliquer cette tension à chaque enroulement.

\begin{center}\TLCouplagesEtoileTriangle\end{center}

\begin{center}
\small
\begin{tabularx}{\linewidth}{>{\bfseries}l X X}
\toprule
Couplage & Tension sur un enroulement & Courants \\
\midrule
Étoile $Y$ & tension simple $V=U/\sqrt3$ & courant d'enroulement égal au courant de ligne.\\
Triangle $\Delta$ & tension composée $U$ & courant d'enroulement $J=I/\sqrt3$.\\
\bottomrule
\end{tabularx}
\end{center}

\TLMAImage[.66\linewidth]{image41.png}
\TLMAImage[.32\linewidth]{image42.jpeg}

\begin{TLMachineMethod}{Choisir le couplage}
\begin{enumerate}
  \item Lire la tension composée du réseau.
  \item Lire les deux tensions de la plaque moteur ; la plus petite est la tension nominale d'un enroulement.
  \item Choisir $Y$ si la tension composée du réseau est égale à la plus grande tension de la plaque.
  \item Choisir $\Delta$ si elle est égale à la plus petite tension de la plaque.
\end{enumerate}
\end{TLMachineMethod}

\begin{TLMachineExercise}{Couplage d'un moteur $400/690\ \mathrm V$}
Déterminer le couplage sur un réseau $230/400\ \mathrm V$, puis sur un réseau $400/690\ \mathrm V$.
\end{TLMachineExercise}

\subsection{Puissances et couple thermique équivalent}

Un récepteur triphasé équilibré est assimilable à trois récepteurs monophasés identiques. Les puissances sont donc données par les relations établies à la section~\ref{sec:masms-1}.

Lorsqu'un moteur délivre un couple cyclique, sa capacité thermique est évaluée par le \textbf{couple thermique équivalent} :
\[
\boxed{C_{th}=\sqrt{\frac{1}{T}\int_0^T C^2(t)\,\mathrm dt}.}
\]
Pour un cycle à couples constants par morceaux :
\[
\boxed{C_{th}=\sqrt{\frac{\sum_i C_i^2t_i}{\sum_i t_i}}.}
\]
La condition de choix est :
\[
\boxed{C_{th}\leq C_N.}
\]

\TLMAImage[.72\linewidth]{image48.png}

\begin{TLMachineWarning}{Dimensionnement thermique}
Le couple maximal admissible pendant un court instant et le couple nominal permanent sont deux caractéristiques différentes. Un cycle doit être vérifié par le couple thermique équivalent et par les surcharges maximales admissibles.
\end{TLMachineWarning}

% ============================================================================
% COURS 13 — MACHINES ASYNCHRONES ET SYNCHRONES
% Partie B : machine asynchrone
% Source : 17-MAS_MS.docx
% ============================================================================

\section{Machine asynchrone}
\label{sec:masms-4}

\subsection{Principe et glissement}

Dans une machine asynchrone, le champ tournant statorique balaie les conducteurs du rotor et y induit des forces électromotrices. Les circuits rotoriques étant fermés, des courants induits apparaissent. L'action du champ statorique sur ces courants crée le couple électromagnétique.

Le rotor ne peut pas atteindre la vitesse du champ : au synchronisme, il ne verrait plus de variation de flux, la f.é.m. rotorique et le couple deviendraient nuls.

On définit le \textbf{glissement} :
\[
\boxed{g=\frac{\Omega_s-\Omega}{\Omega_s}=\frac{N_s-N}{N_s}.}
\]

La fréquence des courants rotoriques est :
\[
\boxed{f_r=g f,}
\qquad
\boxed{\omega_r=g\omega.}
\]

Les principaux domaines de fonctionnement sont :
\begin{center}
\small
\begin{tabularx}{\linewidth}{>{\bfseries}c c X}
\toprule
Glissement & Vitesse & Fonctionnement \\
\midrule
$0<g<1$ & $0<\Omega<\Omega_s$ & moteur.\\
$g<0$ & $\Omega>\Omega_s$ & génératrice ou freinage hypersynchrone.\\
$g>1$ & $\Omega<0$ avec champ direct & freinage par contre-courant.\\
$g=1$ & $\Omega=0$ & démarrage, rotor bloqué.\\
$g=0$ & $\Omega=\Omega_s$ & couple nul.\\
\bottomrule
\end{tabularx}
\end{center}

\begin{TLMachineWarning}{Pertes rotoriques}
Plus le glissement est élevé, plus les pertes Joule rotoriques sont importantes. Le point nominal d'une MAS industrielle se situe généralement à faible glissement.
\end{TLMachineWarning}

\subsection{Schéma monophasé équivalent}

En régime permanent triphasé équilibré, la MAS peut être étudiée par un schéma équivalent monophasé. Les grandeurs rotoriques sont ramenées au stator.

\begin{center}\TLSchemaEquivalentMAS\end{center}

Les éléments du modèle sont :
\begin{itemize}
  \item $R_1$ : résistance d'une phase statorique ;
  \item $X_1=L_1\omega$ : réactance de fuite statorique ;
  \item $R_f$ : résistance modélisant les pertes fer ;
  \item $X_0=L_0\omega$ : réactance de magnétisation ;
  \item $R$ et $X=L\omega$ : résistance et réactance de fuite rotoriques ramenées au stator ;
  \item $R/g$ : résistance fictive traduisant la puissance transmise au rotor.
\end{itemize}

La puissance transmise du stator au rotor vaut :
\[
\boxed{P_{tr}=3\frac{R}{g}{I_1'}^2.}
\]
Les pertes fer statoriques sont modélisées par :
\[
\boxed{P_{fs}=3\frac{V_1^2}{R_f}.}
\]

\TLMAImage[.54\linewidth]{image52.png}

\subsection{Bilan de puissance et rendement}

\begin{center}\TLBilanPuissanceMAS\end{center}

Le bilan de puissance s'écrit :
\[
P_1-P_{Js}-P_{fs}=P_{tr},
\]
\[
\boxed{P_{Jr}=gP_{tr},}
\qquad
\boxed{P_{em}=(1-g)P_{tr}=C_{em}\Omega.}
\]
Après retrait des pertes mécaniques :
\[
P_u=P_{em}-P_{\mathrm{mec}},
\qquad
\boxed{\eta=\frac{P_u}{P_1}.}
\]

La relation $P_{tr}=C_{em}\Omega_s$ est particulièrement utile :
\[
\boxed{C_{em}=\frac{P_{tr}}{\Omega_s}.}
\]

\TLMAImage[.82\linewidth]{image53.png}

\begin{TLMachineExercise}{Bilan nominal d'une MAS}
Une machine porte les indications $230/400\ \mathrm V$, $50\ \mathrm{Hz}$, $15\ \mathrm{kW}$, $1440\ \mathrm{tr\,min^{-1}}$. À charge nominale, le courant de ligne vaut $19{,}4\ \mathrm A$ et la puissance absorbée $11\ \mathrm{kW}$. Les pertes fer et mécaniques sont négligées.
\begin{enumerate}
  \item Déterminer le nombre de paires de pôles et le glissement.
  \item Calculer la puissance réactive et le facteur de puissance.
  \item Déterminer le couple nominal et les pertes Joule rotoriques.
  \item Identifier les paramètres du schéma équivalent à partir d'un essai à vide et d'un essai rotor bloqué.
\end{enumerate}
\end{TLMachineExercise}

\subsection{Couple électromagnétique}

Pour l'étude du couple, les chutes de tension dans $R_1$ et $X_1$ sont souvent négligées. La branche rotorique est alors alimentée sous $V_1$ :
\[
\underline Z_1=\frac{R}{g}+jX,
\qquad
I_1'=\frac{V_1}{\sqrt{(R/g)^2+X^2}}.
\]

En utilisant $P_{tr}=3(R/g){I_1'}^2=C_{em}\Omega_s$ :
\[
\boxed{C_{em}=\frac{3V_1^2}{\Omega_s}\frac{Rg}{R^2+(Xg)^2}.}
\]
Une écriture équivalente est :
\[
C_{em}=\frac{p}{\omega}\,
\frac{3V_1^2(R/g)}{(R/g)^2+X^2}.
\]

\begin{center}\TLCoupleGlissementMAS\end{center}
\TLMAImage[.68\linewidth]{image56.png}

Le couple maximal est obtenu pour :
\[
\boxed{g_M=\frac{R}{X}.}
\]
Sa valeur est indépendante de $R$ :
\[
\boxed{C_M=\frac{3V_1^2}{2\Omega_sX}.}
\]
Au démarrage, $g=1$ :
\[
\boxed{C_D=\frac{3V_1^2}{\Omega_s}\frac{R}{R^2+X^2}.}
\]

Au voisinage du synchronisme, $g\ll g_M$, la caractéristique est presque linéaire :
\[
\boxed{C_{em}\simeq\frac{3V_1^2}{\Omega_sR}\,g.}
\]

\begin{TLMachineMethod}{Étudier une MAS à partir du schéma équivalent}
\begin{enumerate}
  \item Déterminer $p$, $\Omega_s$ et $g$.
  \item Identifier les paramètres par essai à vide et rotor bloqué.
  \item Calculer $P_{tr}$ puis $C_{em}=P_{tr}/\Omega_s$.
  \item Rechercher $g_M$ et $C_M$.
  \item Placer le point de fonctionnement par l'intersection entre couple moteur et couple résistant.
  \item Vérifier les courants, le rendement et la stabilité.
\end{enumerate}
\end{TLMachineMethod}

\begin{TLMachineExercise}{Expression et maximum du couple}
À partir du schéma simplifié comportant $R/g$ et $l\omega$, montrer que :
\[
C=\frac{3pV^2}{\omega}\frac{R/g}{(R/g)^2+(l\omega)^2}.
\]
Déterminer $g_M$, $C_M$, la vitesse correspondante et tracer la caractéristique pour une machine dont la vitesse de synchronisme vaut $1500\ \mathrm{tr\,min^{-1}}$.
\end{TLMachineExercise}

\subsection{Variation de vitesse}

La vitesse mécanique est :
\[
\boxed{\Omega=\frac{\omega}{p}(1-g).}
\]
Elle peut être modifiée en agissant sur :
\begin{itemize}
  \item le glissement, notamment par la tension ;
  \item la fréquence statorique ;
  \item le nombre de paires de pôles.
\end{itemize}

La solution moderne consiste à alimenter la machine par un onduleur et à modifier la fréquence. Pour conserver le flux magnétique approximativement constant, on maintient :
\[
\boxed{\frac{V_1}{f}=\mathrm{constante}.}
\]

\begin{center}\TLCommandeVsurF\end{center}

Sous la fréquence nominale, le flux et le couple maximal restent presque constants. Au-delà de la fréquence nominale, la tension ne peut plus augmenter : le fonctionnement se poursuit à puissance approximativement constante et le couple disponible diminue.

\TLMAImage[.56\linewidth]{image58.png}
\TLMAImage[.56\linewidth]{image59.png}

Dans la zone utile et pour un faible glissement, la caractéristique peut s'écrire sous forme affine :
\[
C_{em}\simeq K(\Omega_s-\Omega).
\]
Cette propriété explique la translation des courbes couple--vitesse lorsque la fréquence varie.

\begin{TLMachineWarning}{Limites de la commande scalaire}
La loi $V/f=\mathrm{cte}$ compense correctement le flux à vitesse moyenne et élevée. À basse fréquence, la chute de tension statorique devient comparable à la tension appliquée : une compensation ou une commande vectorielle est alors nécessaire.
\end{TLMachineWarning}

\begin{TLMachineExercise}{Identification et commande $V/f$}
Une MAS $230/400\ \mathrm V$, $50\ \mathrm{Hz}$, $1440\ \mathrm{tr\,min^{-1}}$ est caractérisée par :
\begin{itemize}
  \item essai à vide : $V=230\ \mathrm V$, $I_0=9{,}5\ \mathrm A$, $P_0=630\ \mathrm W$ ;
  \item essai rotor bloqué : $V_{cc}=50\ \mathrm V$, $I_{cc}=30\ \mathrm A$, $P_{cc}=830\ \mathrm W$.
\end{itemize}
Déterminer les paramètres du modèle, établir l'expression du couple et montrer l'effet d'une commande $V/f$ constante sur les caractéristiques couple--vitesse.
\end{TLMachineExercise}

\subsection{Fonctionnement dans les quatre quadrants}

Une association onduleur--MAS réversible peut fonctionner en moteur ou en génératrice dans les deux sens de rotation. Le signe de la puissance mécanique est donné par :
\[
P_m=C\Omega.
\]

\begin{center}\TLQuadrantsMachine\end{center}

La permutation de deux phases inverse le sens du champ tournant. Pour permettre le freinage régénératif, la chaîne d'énergie doit être réversible jusqu'à la source ; sinon, l'énergie est dissipée dans une résistance de freinage.

% ============================================================================
% COURS 13 — MACHINES ASYNCHRONES ET SYNCHRONES
% Partie C : machine synchrone
% Source : 17-MAS_MS.docx
% ============================================================================

\section{Machine synchrone}
\label{sec:masms-5}

Une machine synchrone est un convertisseur électromécanique réversible. Elle fonctionne en moteur ou en génératrice, appelée alors \textbf{alternateur}. Le rotor crée un champ magnétique par aimants permanents ou par un enroulement d'excitation. En régime établi, il tourne exactement à la vitesse du champ tournant statorique :
\[
\boxed{\Omega=\Omega_s=\frac{\omega}{p}.}
\]

\subsection{Domaines d'emploi}

Les machines synchrones couvrent une très large gamme de puissances :
\begin{itemize}
  \item petites puissances : modélisme, ventilateurs, dispositifs médicaux et micromécanismes ;
  \item puissances moyennes : machines-outils, véhicules électriques, entraînements directs, alternateurs automobiles ;
  \item fortes puissances : traction ferroviaire, propulsion navale et production d'électricité.
\end{itemize}

L'absence de glissement permet une relation directe entre fréquence électrique et vitesse mécanique. Les aimants permanents offrent une forte densité de puissance ; l'excitation bobinée permet d'agir sur le flux et la puissance réactive.

\subsection{Modèle simplifié d'une phase}

Le modèle monophasé équivalent comporte une f.é.m. interne $\underline E$, la résistance statorique $R_s$ et la réactance synchrone $X_s=L_s\omega$.

\begin{center}\TLSchemaEquivalentMS\end{center}

En convention récepteur, pour un moteur :
\[
\boxed{\underline V=\underline E+(R_s+jX_s)\underline I.}
\]
En négligeant la résistance :
\[
\boxed{\underline V=\underline E+jX_s\underline I.}
\]

La f.é.m. est proportionnelle au flux et à la vitesse :
\[
\boxed{E=K_e\Omega_s=p\Phi\Omega_s.}
\]

\begin{TLMachineWarning}{Domaine de validité}
Le modèle à inductance synchrone unique est établi en régime permanent et convient surtout aux machines à pôles lisses. Les machines saillantes demandent généralement deux réactances, directe et transverse.
\end{TLMachineWarning}

\subsection{Diagramme de Behn--Eschenburg}

Le diagramme de Behn--Eschenburg est le diagramme de Fresnel associé à la loi des mailles.

\begin{center}\TLDiagrammeBehnEschenburg\end{center}

Les angles utilisés sont :
\begin{itemize}
  \item $\varphi$ : déphasage entre $\underline I$ et $\underline V$ ;
  \item $\delta$ : angle interne entre $\underline E$ et $\underline V$ ;
  \item $\psi$ : déphasage entre $\underline I$ et $\underline E$.
\end{itemize}

Lorsque $R_s$ est négligée, le vecteur $jX_s\underline I$ est orthogonal au courant. La valeur de $E$, réglée par l'excitation pour une machine bobinée, modifie le facteur de puissance :
\begin{itemize}
  \item machine sous-excitée : comportement généralement inductif ;
  \item machine surexcitée : comportement capacitif possible ;
  \item excitation adaptée : facteur de puissance proche de l'unité.
\end{itemize}

\begin{TLMachineMethod}{Construire un diagramme de Behn--Eschenburg}
\begin{enumerate}
  \item Choisir la grandeur de référence, souvent $\underline E$ ou $\underline V$.
  \item Placer le courant avec le déphasage imposé.
  \item Construire $R_s\underline I$, colinéaire au courant.
  \item Construire $jX_s\underline I$, en quadrature avance.
  \item Fermer le polygone vectoriel avec la loi des mailles.
  \item Lire les angles $\varphi$, $\psi$ et $\delta$ ainsi que les modules recherchés.
\end{enumerate}
\end{TLMachineMethod}

\subsection{Puissance et couple électromagnétiques}

La puissance électrique active absorbée par une machine triphasée équilibrée est :
\[
P=3VI\cos\varphi.
\]
La puissance électromagnétique peut également s'écrire :
\[
\boxed{P_{em}=3EI\cos\psi.}
\]
On en déduit :
\[
\boxed{C_{em}=\frac{P_{em}}{\Omega_s}=3K_eI\cos\psi.}
\]

Si la résistance statorique est négligée et pour une machine non saillante :
\[
\boxed{P_{em}=3\frac{VE}{X_s}\sin\delta,}
\qquad
\boxed{C_{em}=\frac{3VE}{X_s\Omega_s}\sin\delta.}
\]

Le couple est maximal pour $|\delta|=\pi/2$. Au-delà, la machine entre dans une zone instable et risque le \textbf{décrochage}.

L'angle géométrique $\theta$ entre le champ rotorique et le champ statorique est complémentaire de $\psi$ :
\[
\theta+\psi=\frac{\pi}{2},
\qquad
\sin\theta=\cos\psi.
\]
Ainsi :
\[
C_{em}=3K_eI\sin\theta.
\]

\begin{TLMachineWarning}{Décrochage}
Une machine synchrone directement alimentée par un réseau ne peut pas accepter un couple résistant supérieur à son couple maximal. Le rotor décroche du champ tournant et la machine doit être arrêtée puis redémarrée.
\end{TLMachineWarning}

\begin{TLMachineExercise}{Identification d'une MS à aimants permanents}
Une MS possède huit pôles, des enroulements couplés en étoile et un courant nominal $I_N=155\ \mathrm A$. Les essais donnent :
\begin{itemize}
  \item résistance mesurée entre deux phases : $r=0{,}06\ \Omega$ ;
  \item à vide, à $1500\ \mathrm{tr\,min^{-1}}$ : $E=37\ \mathrm V$ par phase ;
  \item en moteur, à $1500\ \mathrm{tr\,min^{-1}}$ : $I=185\ \mathrm A$, $V=49\ \mathrm V$ et $\psi=0$.
\end{itemize}
Déterminer $f$, la résistance d'une phase, la constante $K_e$, l'inductance synchrone puis construire le diagramme vectoriel pour un fonctionnement à $5000\ \mathrm{tr\,min^{-1}}$.
\end{TLMachineExercise}

\subsection{Bilan de puissance}

Les pertes d'une machine synchrone sont :
\begin{itemize}
  \item pertes Joule statoriques : $P_{Js}=3R_sI^2$ ;
  \item pertes fer $P_{fe}$ ;
  \item pertes mécaniques $P_m$ ;
  \item pertes d'excitation $P_e=U_eI_e=R_eI_e^2$ lorsque le rotor est bobiné.
\end{itemize}

\begin{center}\TLBilanPuissanceMS\end{center}

En fonctionnement moteur :
\[
P_{em}=P_a-P_{Js}-P_{fe},
\qquad
P_u=P_{em}-P_m,
\qquad
\boxed{\eta=\frac{P_u}{P_a+P_e}.}
\]
Pour une machine à aimants permanents, $P_e=0$.

\subsection{Variation de vitesse}

La vitesse étant imposée par la fréquence, la variation de vitesse exige une alimentation à fréquence variable. Le contrôle peut agir sur :
\begin{itemize}
  \item le flux ou l'excitation, si le rotor est bobiné ;
  \item l'amplitude du courant statorique ;
  \item l'angle $\psi$ entre courant et f.é.m.
\end{itemize}

Pour une machine à aimants, l'objectif usuel est de maximiser le couple par ampère en imposant $\psi=0$, donc $\cos\psi=1$.

\subsection{Moteur brushless et autopilotage}

Une machine synchrone autopilotée associe :
\begin{itemize}
  \item un onduleur triphasé ;
  \item une mesure ou une estimation de la position rotorique ;
  \item une commande qui synchronise les courants statoriques avec le rotor.
\end{itemize}

\begin{center}\TLChaineBrushless\end{center}

La commande impose une référence d'amplitude $I_{ref}$ et un angle d'autopilotage $\psi_{ref}$. La fonction assurée par l'électronique est comparable à celle du collecteur mécanique d'une machine à courant continu, d'où le nom \textbf{brushless}.

\begin{TLMachineExercise}{Unité de vissage brushless}
Une MS possède $p=3$, $L=7{,}5\ \mathrm{mH}$ et $K_e=0{,}13\ \mathrm{V/(rad\,s^{-1})}$. Elle tourne à $3600\ \mathrm{tr\,min^{-1}}$.
\begin{enumerate}
  \item Calculer $\omega$ et $f$ des grandeurs statoriques.
  \item Montrer que $C=K_cI\cos\psi$ et déterminer $K_c$.
  \item Choisir $\psi$ pour minimiser le courant en mode moteur puis en freinage.
  \item Tracer les diagrammes de Behn--Eschenburg correspondants.
  \item Pour $N=3000\ \mathrm{tr\,min^{-1}}$ et $I=6{,}7\ \mathrm A$, déterminer le facteur de puissance, la tension simple et le couple.
\end{enumerate}
\end{TLMachineExercise}

\section{Synthèse et démarche d'étude}
\label{sec:masms-6}

\begin{TLMachineMethod}{Analyser une machine alternative}
\begin{enumerate}
  \item Lire la plaque signalétique et identifier le couplage.
  \item Déterminer $p$, $\Omega_s$ et les grandeurs triphasées.
  \item Pour une MAS, calculer le glissement et exploiter le schéma équivalent.
  \item Pour une MS, établir le diagramme de Behn--Eschenburg.
  \item Construire le bilan de puissance et calculer le rendement.
  \item Déterminer le couple et le point de fonctionnement avec la charge.
  \item Vérifier les contraintes thermiques et le domaine stable.
  \item Choisir la stratégie de variation de vitesse : $V/f$, commande vectorielle ou autopilotage.
\end{enumerate}
\end{TLMachineMethod}

\begin{center}
\small
\begin{tabularx}{\linewidth}{>{\bfseries}p{3.2cm}X X}
\toprule
Critère & Machine asynchrone & Machine synchrone \\
\midrule
Vitesse & différente de $\Omega_s$ ; dépend du glissement & égale à $\Omega_s$ en régime établi\\
Rotor & cage ou rotor bobiné & aimants ou excitation bobinée\\
Couple & créé par courants rotoriques induits & interaction entre champs statorique et rotorique\\
Commande simple & robuste, loi $V/f$ & nécessite synchronisation avec le rotor\\
Risque principal & échauffement, zone instable & décrochage\\
Atouts & simplicité, coût, robustesse & rendement, précision de vitesse, densité de puissance\\
\bottomrule
\end{tabularx}
\end{center}

\section{Sources}
\label{sec:masms-7}

Ce cours a été élaboré à partir du document pédagogique original et de ressources de collègues de l'UPSTI. Les exercices conservent les références de concours indiquées dans le document source.

% ============================================================================
% COURS 13 — EXERCICES, PARTIE A
% Source : 17-MAS_MS.docx
% ============================================================================

\section{Exercices du chapitre}
\label{sec:masms-8}

\subsection{Banc BALAFRE}

\textit{Source : concours ATS 2019.}

\begin{center}
\TLMAInlineImage[.28\linewidth]{image79.png}\hfill
\TLMAInlineImage[.62\linewidth]{image80.png}
\end{center}

Le banc BALAFRE (BAnc d'essais à LAmes Fluides à haut REynolds) permet d'identifier le comportement dynamique de joints annulaires utilisés dans les turbomachines. Un moteur asynchrone triphasé entraîne le rotor du banc. Pendant l'essai, sa vitesse doit être stabilisée à $6000\ \mathrm{tr\,min^{-1}}$ malgré un couple résistant estimé à $C_{res}=300\ \mathrm{N\,m}$.

Le réseau fournit $230/400\ \mathrm V$ à $50\ \mathrm{Hz}$. La machine est commandée par un variateur fonctionnant à $U_s/f$ constant. Les pertes Joule statoriques, les pertes fer et les pertes mécaniques sont négligées.

\TLMAImage[.45\linewidth]{image81.png}

\begin{enumerate}
  \item À partir de la plaque signalétique, déterminer le couplage, le nombre de paires de pôles, la vitesse de synchronisme, le glissement nominal et le couple utile nominal.
  \item Le modèle équivalent d'une phase comporte une inductance de magnétisation $L_0$, une inductance de fuite totale $L_c$ et une résistance rotorique ramenée au stator $R/g$. Établir :
  \[
  C_u=\frac{3pU_s^2}{\omega}\,
  \frac{R/g}{(R/g)^2+(L_c\omega)^2}.
  \]
  \item Identifier sur la caractéristique couple--vitesse le démarrage, le point nominal, le synchronisme et la zone instable.
  \item Déterminer le glissement $g_M$ correspondant au couple maximal et exprimer $C_M$.
  \item Le constructeur indique $C_M/C_N=3{,}5$. Déterminer $L_c$, puis estimer $R$ à partir du point nominal.
  \item À l'aide des courbes fournies pour des fréquences comprises entre $90$ et $110\ \mathrm{Hz}$, déterminer la fréquence à imposer pour maintenir $6000\ \mathrm{tr\,min^{-1}}$ sous $300\ \mathrm{N\,m}$.
  \item Vérifier que la durée d'accélération imposée par le cahier des charges reste inférieure à $5\ \mathrm s$.
\end{enumerate}

\subsection{Panneaux publicitaires déroulants}

\textit{Source : concours ATS 2011.}

\begin{minipage}{.73\linewidth}
Le panneau déroulant de type Senior utilise deux moteurs asynchrones commandés par des variateurs scalaires. Les affiches défilent à $1\ \mathrm{m\,s^{-1}}$ avec des rampes d'accélération et de décélération d'une seconde.

En régime permanent, le couple moteur nécessaire vaut $C_m=0{,}28\ \mathrm{N\,m}$ pour une vitesse $\Omega_m=139\ \mathrm{rad\,s^{-1}}$. Le motoréducteur possède les caractéristiques nominales suivantes sous $50\ \mathrm{Hz}$ :
\[
P_N=120\ \mathrm W,\quad N_N=1300\ \mathrm{tr\,min^{-1}},\quad
C_N=0{,}88\ \mathrm{N\,m},\quad C_{max}/C_N=1{,}9.
\]
\end{minipage}\hfill
\begin{minipage}{.23\linewidth}
\TLMAInlineImage[\linewidth]{image86.jpeg}
\end{minipage}

Le couple est modélisé par :
\[
C_m=\frac{2C_{max}}{g_0/g+g/g_0}.
\]

\begin{enumerate}
  \item Calculer la puissance mécanique nécessaire et montrer que le moteur est surdimensionné.
  \item Déterminer $p$, $\Omega_s$, $N_s$ et le glissement nominal.
  \item Calculer $g_0$, puis tracer $C_m(g)$ pour $0<g<1$ en repérant la zone stable.
  \item Pour le couple réel $C_{maff}=0{,}30\ \mathrm{N\,m}$, utiliser l'approximation de faible glissement pour déterminer $g_{aff}$.
  \item Montrer que, dans cette zone, $C_m=\lambda(\Omega_s-\Omega_m)$ et préciser $\lambda$.
  \item La commande est réalisée à $U/f$ constant et l'on prendra $\lambda=0{,}0453\ \mathrm{N\,m\,s\,rad^{-1}}$. Déterminer la vitesse de synchronisme puis la fréquence à programmer pour obtenir $\Omega_m=139\ \mathrm{rad\,s^{-1}}$.
\end{enumerate}

\subsection{Entraînement de la broche d'un centre d'usinage}

\textit{Source : concours CCP TSI 2002.}

\begin{minipage}{.72\linewidth}
Le centre d'usinage horizontal étudié possède une broche entraînée par une machine asynchrone triphasée et un onduleur de tension. Les caractéristiques nominales sous $50\ \mathrm{Hz}$ sont :
\[
P_N=5{,}35\ \mathrm{kW},\quad C_N=35\ \mathrm{N\,m},\quad
N_N=1460\ \mathrm{tr\,min^{-1}},\quad I_N=9{,}5\ \mathrm A,
\]
avec $C_D/C_N=2{,}4$ et $C_{max}=89{,}3\ \mathrm{N\,m}$.
\end{minipage}\hfill
\begin{minipage}{.24\linewidth}
\TLMAInlineImage[\linewidth]{image88.jpeg}
\end{minipage}

Toutes les pertes sont négligées sauf les pertes Joule rotoriques. Le moteur fonctionne à couple maximal constant jusqu'à $1500\ \mathrm{tr\,min^{-1}}$, puis à puissance utile maximale constante jusqu'à $6000\ \mathrm{tr\,min^{-1}}$.

\begin{enumerate}
  \item Déterminer $p$, $N_s$, $g_N$, $P_{tr}$, les pertes Joule rotoriques, le rendement et le facteur de puissance au point nominal.
  \item Tracer l'enveloppe couple--vitesse autorisée du moteur.
  \item Reporter sur le même graphe les couples résistants ramenés au moteur pour les deux rapports de boîte :
  \[
  C_{r1}=12+0{,}15\Omega,
  \qquad
  C_{r2}=1+0{,}0125\Omega.
  \]
  \item Déterminer les points de fonctionnement et les couples accélérateurs disponibles.
  \item L'onduleur est commandé à $U/f$ constant et délivre $400\ \mathrm V$ entre phases à $50\ \mathrm{Hz}$. Tracer le coefficient de modulation en fonction de la fréquence jusqu'à $80\ \mathrm{Hz}$.
  \item À partir du schéma équivalent simplifié, établir l'expression de $C_{em}$, calculer $C_{max}$ à $50\ \mathrm{Hz}$ et rechercher la fréquence donnant le couple accélérateur maximal au démarrage.
\end{enumerate}

\subsection{Maison hantée}

\textit{Source : concours ATS 2009.}

\begin{center}
\TLMAInlineImage[.28\linewidth]{image92.png}\hfill
\TLMAInlineImage[.62\linewidth]{image93.jpeg}
\end{center}

Seize véhicules circulent sur une piste guidée. Chaque véhicule est entraîné par un moteur asynchrone triphasé associé à un variateur. La vitesse de translation doit rester comprise entre $1$ et $1{,}3\ \mathrm{m\,s^{-1}}$, avec une variation maximale tolérée de $10\,\%$. Le moteur SEW DV100M4, connecté en triangle, possède sous $50\ \mathrm{Hz}$ :
\[
P_N=2{,}2\ \mathrm{kW},\quad C_N=15\ \mathrm{N\,m},\quad
N_N=1410\ \mathrm{tr\,min^{-1}},\quad C_{max}=40\ \mathrm{N\,m}.
\]

Le couple électromagnétique est modélisé par :
\[
C_{em}=\frac{3U^2}{\Omega_s}\frac{R_2g}{R_2^2+(X_2g)^2}.
\]

\begin{enumerate}
  \item Déterminer $g_M$ et $C_M$, puis écrire la relation réduite :
  \[
  C_{em}=2C_M\frac{x}{1+x^2},\qquad x=\frac{g}{g_M}.
  \]
  \item Déterminer $p$ et le glissement nominal. Utiliser les données constructeur pour calculer $g_M$.
  \item Linéariser la caractéristique autour du synchronisme et exprimer $C_{em}$ en fonction de $f$, $p$ et $\Omega_m$.
  \item Déterminer la fréquence qui assure $1{,}3\ \mathrm{m\,s^{-1}}$ pour un couple moteur de $10\ \mathrm{N\,m}$.
  \item Étudier la réversibilité du convertisseur et préciser le trajet de l'énergie lorsque le couple devient négatif.
  \item Calculer la vitesse lorsque le couple passe de $10$ à $-4\ \mathrm{N\,m}$ et conclure sur la nécessité d'une régulation de vitesse.
\end{enumerate}

% ============================================================================
% COURS 13 — EXERCICES, PARTIE B
% Source : 17-MAS_MS.docx
% ============================================================================

\subsection{Télésiège débrayable six places}

\textit{Source : concours CCP TSI 2016.}

\TLMAImage[.70\linewidth]{image97.jpeg}

Le télésiège « Biollay » possède deux moteurs de secours asynchrones triphasés de $75\ \mathrm{kW}$. Ils entraînent la poulie motrice de rayon $R_p=2{,}45\ \mathrm m$ par un premier réducteur de rapport $r_1=32{,}7$, puis un engrenage comportant une couronne de $220$ dents et un pignon de $16$ dents.

Un moteur possède les caractéristiques :
\[
U=400\ \mathrm V,\quad I=133\ \mathrm A,\quad f=50\ \mathrm{Hz},\quad
N_N=2978\ \mathrm{tr\,min^{-1}},\quad \eta=93{,}8\,\%,\quad \cos\varphi=0{,}87.
\]

\begin{enumerate}
  \item Exprimer le couple d'un moteur en fonction des tensions du câble $T_A$, $T_B$, du rayon de la poulie et des rapports de transmission. Calculer ce couple au début puis à la fin de l'évacuation.
  \item Déterminer $N_s$, $p$ et le glissement nominal.
  \item Le modèle rotorique est défini par $V_1=230\ \mathrm V$, $R_2'=10{,}5\ \mathrm{m\Omega}$ et $X_2'=0{,}23\ \Omega$. Exprimer le courant $I_2'$, la puissance transmise et le couple électromagnétique.
  \item Déterminer le glissement $g_M$ du couple maximal.
  \item En mode dégradé, un seul moteur doit fournir $420\ \mathrm{N\,m}$ au début et $120\ \mathrm{N\,m}$ à la fin. Calculer les glissements et les vitesses du câble correspondantes.
  \item Vérifier que la vitesse reste entre $0{,}8$ et $1{,}8\ \mathrm{m\,s^{-1}}$ et valider le choix du moteur.
\end{enumerate}

\subsection{Véhicule autoguidé hospitalier}

\textit{Source : concours ATS 2014.}

\TLMAImage[.68\linewidth]{image100.png}

Le véhicule autonome à guidage laser est entraîné par un moteur-roue asynchrone. La roue a un rayon $r=105\ \mathrm{mm}$ et le réducteur possède un rapport $r_1$. La vitesse de translation doit pouvoir varier entre $0{,}1$ et $1{,}2\ \mathrm{m\,s^{-1}}$.

Le moteur, couplé en étoile, possède les caractéristiques nominales :
\[
f_s=110\ \mathrm{Hz},\quad U_n=14{,}5\ \mathrm V,\quad
N_n=3100\ \mathrm{tr\,min^{-1}},\quad P_u=850\ \mathrm W,
\]
\[
\cos\varphi=0{,}79,\qquad \eta_n=0{,}76.
\]
La source continue du variateur vaut $24\ \mathrm V$.

\begin{enumerate}
  \item Établir la relation entre $N_s$, $f$ et $p$, puis déterminer $p$ et $N_s$.
  \item Calculer le glissement nominal et le couple nominal.
  \item Établir la relation entre la vitesse du véhicule, la fréquence et le glissement.
  \item Pour $g=6\,\%$, mettre la loi entrée--sortie sous la forme $V_A=Kf$ et déterminer la plage de fréquence nécessaire.
  \item La consigne de vitesse du variateur est codée sur un mot signé de 16 bits avec un pas de $0{,}25\ \mathrm{tr\,min^{-1}}$. Déterminer le mot de commande correspondant à $2620\ \mathrm{tr\,min^{-1}}$.
  \item La commande maintient $U/f$ constant. Déterminer la plage de tension efficace nécessaire et vérifier sa compatibilité avec la batterie.
\end{enumerate}

\subsection{Dépose bagage automatique}

\textit{Source : concours Centrale--Supélec TSI 2018.}

Le basculeur d'un système de dépose bagage automatique est entraîné par une MAS triphasée, un variateur et un réducteur de rapport $1/107{,}7$. La plaque indique :
\[
P_u=730\ \mathrm W,\quad 230/400\ \mathrm V,\quad f=50\ \mathrm{Hz},\quad
N_N=1430\ \mathrm{tr\,min^{-1}},\quad C_N=4{,}9\ \mathrm{N\,m},
\]
avec $C_D/C_N=1{,}16$ et $C_{max}/C_N=2{,}6$.

Le modèle simplifié comporte $R/g$ en série avec $X=L\omega$.

\begin{enumerate}
  \item Déterminer la tension d'un enroulement, le couplage, $p$ et $N_s$.
  \item Montrer que :
  \[
  C_m=\frac{3V^2}{\Omega_s}\frac{R/g}{(R/g)^2+X^2}.
  \]
  \item Déterminer $g_{max}$ et $C_{max}$, puis calculer $L$ à partir du rapport $C_{max}/C_N$.
  \item Utiliser le couple de démarrage pour déterminer la plus petite valeur admissible de $R$.
  \item Le couple résistant est assimilé à $C_r=K_r\omega_{mot}$. Identifier $K_r$ à partir du chronogramme fourni.
  \item Dans la partie utile, utiliser $C_m=(p/\omega)(3V^2/R)g$ et $V/f=\mathrm{cte}$ pour déterminer la fréquence imposant $135{,}3\ \mathrm{rad\,s^{-1}}$.
  \item Comparer ce résultat avec la résolution numérique obtenue par minimisation de $C_m-C_r$.
\end{enumerate}

\TLMAImage[.70\linewidth]{image109.png}

\subsection{Débourreuse de noyaux de fonderie}

\textit{Source : concours Centrale--Supélec TSI 2009.}

Deux moteurs asynchrones à balourds transmettent des vibrations à des culasses en aluminium afin de désagréger les noyaux de sable. La masse mise en mouvement vaut $4940\ \mathrm{kg}$ et l'accélération minimale imposée est $16\ \mathrm{m\,s^{-2}}$ entre $20$ et $25\ \mathrm{Hz}$.

Chaque moteur possède :
\[
P_N=7{,}75\ \mathrm{kW},\quad I_N=13\ \mathrm A,\quad 230/400\ \mathrm V,\quad
f=50\ \mathrm{Hz},\quad N_N=1440\ \mathrm{tr\,min^{-1}}.
\]
Le moment d'inertie ramené sur l'arbre vaut $J=0{,}33\ \mathrm{kg\,m^2}$ et le frottement visqueux $f_v=5\times10^{-2}\ \mathrm{N\,m\,s}$.

\begin{enumerate}
  \item Justifier le choix de deux moteurs à balourds et calculer l'action dynamique demandée à chacun.
  \item Vérifier le choix d'un variateur de $20\ \mathrm{kW}$.
  \item Après coupure de l'alimentation, établir et résoudre l'équation $J\dot\Omega+f_v\Omega=0$. Calculer le temps pour atteindre $1\,\%$ de la vitesse initiale et conclure sur la nécessité d'un freinage.
  \item Déterminer le nombre de paires de pôles.
  \item À partir du modèle comportant $L$, $\ell_{fr}$ et $R_r/g$, établir $P_{tr}$ puis le couple électromagnétique.
  \item Exploiter l'essai au synchronisme, $V=230\ \mathrm V$, $Q=800\ \mathrm{var}$, pour déterminer $L$.
  \item Exploiter l'essai rotor bloqué, $V_{cc}=76\ \mathrm V$, $I_{cc}=13\ \mathrm A$, $P_{cc}=350\ \mathrm W$, pour déterminer $R_r$ et $\ell_{fr}$.
  \item Tracer $C(g)$ pour $-1\leq g\leq1$ et identifier les zones moteur et frein.
\end{enumerate}

\subsection{Tramway de Strasbourg}

\textit{Source : concours troisième année ENS Cachan 2002.}

\TLMAImage[.42\linewidth]{image117.jpeg}
\TLMAImage[.75\linewidth]{image118.png}

Une rame possède douze moteurs asynchrones de traction. Un moteur, couplé en étoile, est caractérisé par :
\[
U_N=585\ \mathrm V,\quad f_N=88\ \mathrm{Hz},\quad I_N=35{,}4\ \mathrm A,\quad
\cos\varphi_N=0{,}732,
\]
\[
N_s=2640\ \mathrm{tr\,min^{-1}},\qquad N_N=2610\ \mathrm{tr\,min^{-1}}.
\]
Le modèle simplifié utilise $R=0{,}138\ \Omega$, $L_t=2{,}38\ \mathrm{mH}$ et $L_M=26{,}55\ \mathrm{mH}$.

\begin{enumerate}
  \item Déterminer $p$, $g_N$, la puissance absorbée, $P_{tr}$, le couple nominal, les pertes Joule rotoriques et la puissance utile.
  \item Établir le couple en fonction de $V$, $L_t$, $R$, $g$, $p$ et $\omega$.
  \item Déterminer le glissement et la valeur du couple maximal.
  \item Tracer $C(N)$ pour $0<N<2N_s$ et préciser les parties stables, moteur et génératrice.
  \item Montrer qu'à faible glissement $C\simeq Kg$ avec $K=8972\ \mathrm{N\,m}$, puis calculer les vitesses pour $C=90\ \mathrm{N\,m}$ et $C=-90\ \mathrm{N\,m}$.
  \item Expliquer l'inversion du sens de rotation et compléter les diagrammes des quatre quadrants mécaniques.
\end{enumerate}

\subsection{Table de radiologie D²RS}

\textit{Source : concours ATS 2018.}

La table radiologique est alimentée par le réseau triphasé $400\ \mathrm V$, $50\ \mathrm{Hz}$. Son mouvement de chariotage est perturbé par des instabilités à basse vitesse.

Le modèle par phase comporte $R_1$, $L_1$, l'inductance de magnétisation $L_m$, puis la branche rotorique $R_2/g$ et $L_2$. Les essais conduisent aux valeurs retenues :
\[
L_m=0{,}382\ \mathrm H,\qquad L_2=29{,}4\ \mathrm{mH},\qquad R_2=9{,}08\ \Omega.
\]

\begin{enumerate}
  \item À partir de la plaque signalétique, déterminer couplage, $p$, glissement, puissance absorbée, pertes et rendement nominal.
  \item Exploiter l'essai à vide sous $230\ \mathrm V$, $Q_0=1322\ \mathrm{var}$ pour retrouver $L_m$.
  \item Exploiter l'essai rotor bloqué sous $15\ \mathrm V$, $I=2{,}49\ \mathrm A$, $S=245\ \mathrm{VA}$ et $f_p=0{,}69$ pour retrouver $R_2$ et $L_2$.
  \item Établir l'expression de $C_{em}$, déterminer $g_{max}$ et $C_{max}$.
  \item Justifier la commande scalaire $V_1/f$ constante à partir du courant magnétisant.
  \item Calculer la chute de tension statorique avec $R_1=0{,}85\ \Omega$ et $L_1=3{,}24\ \mathrm{mH}$.
  \item Expliquer les instabilités à basse fréquence et proposer une commande vectorielle avec compensation des chutes statoriques.
\end{enumerate}

\subsection{Éolienne à génératrice synchrone}

\textit{Source : CAPET SII 2010.}

Une éolienne doit fournir au moins $5\ \mathrm{kW}$ au logement d'un camping. À $240\ \mathrm{tr\,min^{-1}}$, la génératrice synchrone débite $I=8{,}38\ \mathrm A$ avec un déphasage $\varphi=18^\circ$. La f.é.m. interne vaut $E=311\ \mathrm V$ à $f=48\ \mathrm{Hz}$. La puissance mécanique récupérée est $6{,}1\ \mathrm{kW}$.

\begin{enumerate}
  \item En convention générateur, écrire $\underline V_1$ en fonction de $\underline E$, $\underline I$, $r_1$, $L_1$ et $\omega$.
  \item Construire le diagramme de Behn--Eschenburg et déterminer graphiquement $V_1$.
  \item Calculer la puissance active triphasée fournie au réseau.
  \item Vérifier la couverture du besoin de $5\ \mathrm{kW}$ et calculer le rendement de la génératrice.
\end{enumerate}


\printindex
